\documentclass[11pt,a4paper]{article}
\usepackage[utf8]{inputenc}
\usepackage[T1]{fontenc}
\usepackage[spanish]{babel}
\usepackage{geometry}
\usepackage{amsmath,amssymb}
\usepackage{booktabs}
\usepackage{array}
\usepackage{float}
\usepackage{listings}
\usepackage{tikz}
\usetikzlibrary{babel,arrows.meta,positioning}
\usepackage{xcolor}
\usepackage{hyperref}

\geometry{margin=2.3cm}

\lstset{
  basicstyle=\ttfamily\small,
  frame=single,
  breaklines=true,
  backgroundcolor=\color{gray!8},
  keywordstyle=\color{blue!70!black}\bfseries,
  commentstyle=\color{green!45!black},
  showstringspaces=false
}

\title{\textbf{Informe Técnico — Trabajo Práctico 3}\\
Diseño de Gramática LL(1), Reconocedor Descendente y Evaluador para Expresiones Trigonométricas}
\author{Asignatura: Compiladores — Ingeniería de Sistemas}
\date{\today}

\begin{document}
\maketitle

\section{Introducción y Objetivos}
El análisis sintáctico descendente constituye una de las técnicas fundamentales en la construcción de compiladores e intérpretes, permitiendo verificar la estructura gramatical de una entrada y, de manera simultánea, realizar su evaluación semántica mediante traducción dirigida por sintaxis.

El presente trabajo implementa un módulo completo en C++17 capaz de analizar léxica y sintácticamente expresiones aritméticas y trigonométricas, detectar variables simbólicas, solicitar sus valores numéricos al usuario y calcular el resultado con precisión \texttt{double}, reutilizando la arquitectura limpia y el motor léxico AFD desarrollados previamente en \texttt{scanner\_c}.

\subsection*{Objetivos}
\begin{itemize}
    \item Diseñar una gramática libre de contexto \textbf{LL(1)} sin ambigüedad ni recursividad izquierda para expresiones aritméticas y trigonométricas.
    \item Reutilizar el analizador léxico AFD de \texttt{scanner\_c} mediante una especificación de lenguaje dedicada (\texttt{EspecificacionTrigonometrica}).
    \item Implementar un reconocedor descendente recursivo que valide todas las reglas léxicas, sintácticas y matemáticas posibles.
    \item Evaluar directamente las expresiones en precisión \texttt{double} sustituyendo variables simbólicas detectadas.
\end{itemize}

\section{Integración con el Analizador Léxico (\texttt{scanner\_c})}
El proyecto \texttt{scanner\_c} separa el motor AFD (\texttt{EscanerAFD}) de la definición concreta del lenguaje a través de la interfaz \texttt{IEspecificacionLenguaje}. Gracias a este diseño, se reutilizó el motor AFD sin modificar su lógica de escaneo, añadiendo únicamente:
\begin{itemize}
    \item Las categorías \texttt{FN\_SIN}, \texttt{FN\_COS} y \texttt{FN\_TAN} en \texttt{TipoComponente}.
    \item La clase \texttt{EspecificacionTrigonometrica}, que registra como palabras reservadas \texttt{sin}, \texttt{sen} (alias en español presente en los ejemplos del enunciado), \texttt{cos} y \texttt{tan}.
    \item La validación \texttt{esTokenPermitido()}, que filtra tokens propios de C (como \texttt{;}, \texttt{=}, cadenas o literales hexadecimales) no válidos en expresiones trigonométricas.
\end{itemize}

\begin{table}[h!]
\centering
\small
\begin{tabular}{lll}
\toprule
\textbf{Lexema} & \textbf{Token (\texttt{TipoComponente})} & \textbf{Rol en la gramática} \\
\midrule
\texttt{sin}, \texttt{sen}, \texttt{cos}, \texttt{tan} & \texttt{FN\_SIN}, \texttt{FN\_COS}, \texttt{FN\_TAN} & \texttt{FUNC} \\
\texttt{+}, \texttt{-}, \texttt{*}, \texttt{/}, \texttt{\^{}} & \texttt{PLUS}, \texttt{MINUS}, \texttt{MUL\_OP}, \texttt{DIV\_OP}, \texttt{XOR\_OP} & Operadores / potencia \\
\texttt{(}, \texttt{)} & \texttt{LPAR}, \texttt{RPAR} & Agrupación y argumentos \\
Reales (\texttt{0}, \texttt{3.14}, \texttt{.5}, \texttt{3.}, \texttt{1.5e-3}) & \texttt{INT\_NUM}, \texttt{FLOAT\_NUM} & \texttt{NUM} \\
Identificadores (\texttt{x}, \texttt{y}, \texttt{z}) & \texttt{ID} & \texttt{VAR} (variable simbólica) \\
\bottomrule
\end{tabular}
\caption{Tabla de tokens del lenguaje de expresiones trigonométricas.}
\end{table}

\section{Diseño de la Gramática LL(1) con Atributos}
Para respetar la precedencia matemática estándar —paréntesis y funciones primero, luego potencia (\texttt{\^{}}) con asociatividad derecha, luego signo unario (\texttt{+}/\texttt{-}), multiplicación/división y finalmente suma/resta— y eliminar la recursividad izquierda, se diseñó una \textbf{Gramática con Atributos} (atributos sintéticos y heredados) siguiendo la notación formal con subíndices ($p, q, r, s, t$):

\begin{itemize}
    \item \textbf{Atributos Sintéticos ($\uparrow$):} se calculan de \emph{abajo hacia arriba} en el árbol de derivación, y dentro de cada producción fluyen de \emph{derecha a izquierda} (hacia el no terminal de la cabeza).
    \item \textbf{Atributos Heredados ($\downarrow$):} se mueven de \emph{izquierda a derecha} en el mismo nivel del árbol de derivación, o de \emph{arriba hacia abajo} en el árbol. En los no terminales recursivos a derecha ($\langle E'\rangle$, $\langle T'\rangle$, $\langle P'\rangle$), el par $_{p\downarrow,\,q\uparrow}$ indica que $p$ es el atributo heredado (acumulado recibido) y $q$ es el atributo sintético (resultado final devuelto).
\end{itemize}

\subsection{Gramática con Atributos (Reglas Semánticas)}
Cada producción asocia ecuaciones semánticas sobre los atributos numéricos de sus símbolos, sin depender de acciones imperativas ni conversión de cadenas en las reglas:

\begin{table}[H]
\centering
\scriptsize
\begin{tabular}{rll}
\toprule
\textbf{\#} & \textbf{Producción con atributos} & \textbf{Ecuaciones semánticas} \\
\midrule
1  & $\langle E\rangle_p \rightarrow \langle T\rangle_q \, \langle E'\rangle_{t\downarrow,\,r\uparrow}$ & $t = q \,;\; p = r$ \\
2  & $\langle E'\rangle_{p\downarrow,\,q\uparrow} \rightarrow + \, \langle T\rangle_r \, \langle E'\rangle_{s\downarrow,\,t\uparrow}$ & $s = p + r \,;\; q = t$ \\
3  & $\langle E'\rangle_{p\downarrow,\,q\uparrow} \rightarrow - \, \langle T\rangle_r \, \langle E'\rangle_{s\downarrow,\,t\uparrow}$ & $s = p - r \,;\; q = t$ \\
4  & $\langle E'\rangle_{p\downarrow,\,q\uparrow} \rightarrow \lambda$ & $q = p$ \\
\midrule
5  & $\langle T\rangle_p \rightarrow \langle U\rangle_q \, \langle T'\rangle_{t\downarrow,\,r\uparrow}$ & $t = q \,;\; p = r$ \\
6  & $\langle T'\rangle_{p\downarrow,\,q\uparrow} \rightarrow * \, \langle U\rangle_r \, \langle T'\rangle_{s\downarrow,\,t\uparrow}$ & $s = p * r \,;\; q = t$ \\
7  & $\langle T'\rangle_{p\downarrow,\,q\uparrow} \rightarrow / \, \langle U\rangle_r \, \langle T'\rangle_{s\downarrow,\,t\uparrow}$ & $s = p / r \,;\; q = t \quad (r \neq 0)$ \\
8  & $\langle T'\rangle_{p\downarrow,\,q\uparrow} \rightarrow \lambda$ & $q = p$ \\
\midrule
9  & $\langle U\rangle_p \rightarrow + \, \langle U\rangle_q$ & $p = q$ \\
10 & $\langle U\rangle_p \rightarrow - \, \langle U\rangle_q$ & $p = -q$ \\
11 & $\langle U\rangle_p \rightarrow \langle P\rangle_q$ & $p = q$ \\
\midrule
12 & $\langle P\rangle_p \rightarrow \langle F\rangle_q \, \langle P'\rangle_{t\downarrow,\,r\uparrow}$ & $t = q \,;\; p = r$ \\
13 & $\langle P'\rangle_{p\downarrow,\,q\uparrow} \rightarrow \texttt{\^{}} \, \langle U\rangle_r$ & $q = p^{\,r}$ \\
14 & $\langle P'\rangle_{p\downarrow,\,q\uparrow} \rightarrow \lambda$ & $q = p$ \\
\midrule
15 & $\langle F\rangle_p \rightarrow ( \, \langle E\rangle_q \, )$ & $p = q$ \\
16 & $\langle F\rangle_p \rightarrow \text{FUNC} \, ( \, \langle E\rangle_q \, )$ & $p = \text{FUNC}(q)$ \\
17 & $\langle F\rangle_p \rightarrow \text{NUM}_q$ & $p = q$ \\
18 & $\langle F\rangle_p \rightarrow \text{VAR}_q$ & $p = q$ \\
\bottomrule
\end{tabular}
\caption{Gramática con atributos sintéticos ($\uparrow$) y heredados ($\downarrow$).}
\end{table}

\subsection{Gramática Subyacente (Sin Atributos)}
Al omitir los subíndices de atributos se obtiene la gramática libre de contexto subyacente utilizada para calcular anulables, conjuntos \text{FIRST}, \text{FOLLOW} y \text{SELECT}:

\begin{align*}
E  &\rightarrow T \, E' \\
E' &\rightarrow + \, T \, E' \;\mid\; - \, T \, E' \;\mid\; \varepsilon \\
T  &\rightarrow U \, T' \\
T' &\rightarrow * \, U \, T' \;\mid\; / \, U \, T' \;\mid\; \varepsilon \\
U  &\rightarrow + \, U \;\mid\; - \, U \;\mid\; P \\
P  &\rightarrow F \, P' \\
P' &\rightarrow \texttt{\^{}} \, U \;\mid\; \varepsilon \\
F  &\rightarrow ( \, E \, ) \;\mid\; \text{FUNC} \, ( \, E \, ) \;\mid\; \text{NUM} \;\mid\; \text{VAR}
\end{align*}

\noindent Nótese que ubicar el signo unario en $U$ y hacer $P' \rightarrow \texttt{\^{}}\,U \mid \varepsilon$ logra simultáneamente que \texttt{-2\^{}2} se evalúe como $-(2^2) = -4$ y que \texttt{2\^{}-2} sea válido y evalúe a $2^{-2} = 0.25$, además de garantizar asociatividad derecha en \texttt{2\^{}3\^{}2} $= 2^{(3^2)} = 512$.

\subsection{Conjuntos FIRST (Primeros), FOLLOW (Siguientes) y Anulables}
Antes de calcular \text{FIRST} y \text{FOLLOW}, se identifican los símbolos y producciones anulables (que derivan $\varepsilon$):
\[
\text{Nonul} = \{E', T', P'\}, \qquad \text{Panul} = \{4, 8, 14\}
\]
\begin{table}[H]
\centering
\small
\begin{tabular}{lll}
\toprule
\textbf{No terminal} & \textbf{PRIMEROS ($\text{FIRST}$)} & \textbf{SIGUIENTES ($\text{FOLLOW}$)} \\
\midrule
$E$  & $\{+, -, (, \text{FUNC}, \text{NUM}, \text{VAR}\}$ & $\{\$, )\}$ \\
$E'$ & $\{+, -, \varepsilon\}$ & $\{\$, )\}$ \\
$T$  & $\{+, -, (, \text{FUNC}, \text{NUM}, \text{VAR}\}$ & $\{+, -, \$, )\}$ \\
$T'$ & $\{*, /, \varepsilon\}$ & $\{+, -, \$, )\}$ \\
$U$  & $\{+, -, (, \text{FUNC}, \text{NUM}, \text{VAR}\}$ & $\{*, /, +, -, \$, )\}$ \\
$P$  & $\{(, \text{FUNC}, \text{NUM}, \text{VAR}\}$ & $\{*, /, +, -, \$, )\}$ \\
$P'$ & $\{\texttt{\^{}}, \varepsilon\}$ & $\{*, /, +, -, \$, )\}$ \\
$F$  & $\{(, \text{FUNC}, \text{NUM}, \text{VAR}\}$ & $\{\texttt{\^{}}, *, /, +, -, \$, )\}$ \\
\bottomrule
\end{tabular}
\caption{Conjuntos PRIMEROS y SIGUIENTES de cada no terminal.}
\end{table}

\subsection{Conjuntos PRIMEROS y de SELECCIÓN por Producción (Prueba LL(1))}
Para cada producción $A \rightarrow \alpha$, el conjunto de selección se define como:
\[
\text{SELECCIÓN}(A \rightarrow \alpha) =
\begin{cases}
\text{FIRST}(\alpha) & \text{si } \varepsilon \notin \text{FIRST}(\alpha) \\
\big(\text{FIRST}(\alpha) \setminus \{\varepsilon\}\big) \cup \text{FOLLOW}(A) & \text{si } \varepsilon \in \text{FIRST}(\alpha)
\end{cases}
\]

\begin{table}[H]
\centering
\scriptsize
\begin{tabular}{rlll}
\toprule
\textbf{\#} & \textbf{Producción ($A \rightarrow \alpha$)} & \textbf{PRIMEROS ($\text{FIRST}(\alpha)$)} & \textbf{SELECCIÓN ($\text{SELECT}$)} \\
\midrule
1  & $E \rightarrow T\,E'$            & $\{+, -, (, \text{FUNC}, \text{NUM}, \text{VAR}\}$ & $\{+, -, (, \text{FUNC}, \text{NUM}, \text{VAR}\}$ \\
2  & $E' \rightarrow +\,T\,E'$        & $\{+\}$                                            & $\{+\}$ \\
3  & $E' \rightarrow -\,T\,E'$        & $\{-\}$                                            & $\{-\}$ \\
4  & $E' \rightarrow \varepsilon$     & $\{\varepsilon\}$                                  & $\{\$, )\}$ \\
\midrule
5  & $T \rightarrow U\,T'$            & $\{+, -, (, \text{FUNC}, \text{NUM}, \text{VAR}\}$ & $\{+, -, (, \text{FUNC}, \text{NUM}, \text{VAR}\}$ \\
6  & $T' \rightarrow *\,U\,T'$        & $\{*\}$                                            & $\{*\}$ \\
7  & $T' \rightarrow /\,U\,T'$        & $\{/\}$                                            & $\{/\}$ \\
8  & $T' \rightarrow \varepsilon$     & $\{\varepsilon\}$                                  & $\{+, -, \$, )\}$ \\
\midrule
9  & $U \rightarrow +\,U$             & $\{+\}$                                            & $\{+\}$ \\
10 & $U \rightarrow -\,U$             & $\{-\}$                                            & $\{-\}$ \\
11 & $U \rightarrow P$                & $\{(, \text{FUNC}, \text{NUM}, \text{VAR}\}$       & $\{(, \text{FUNC}, \text{NUM}, \text{VAR}\}$ \\
\midrule
12 & $P \rightarrow F\,P'$            & $\{(, \text{FUNC}, \text{NUM}, \text{VAR}\}$       & $\{(, \text{FUNC}, \text{NUM}, \text{VAR}\}$ \\
13 & $P' \rightarrow \texttt{\^{}}\,U$ & $\{\texttt{\^{}}\}$                               & $\{\texttt{\^{}}\}$ \\
14 & $P' \rightarrow \varepsilon$     & $\{\varepsilon\}$                                  & $\{*, /, +, -, \$, )\}$ \\
\midrule
15 & $F \rightarrow (\,E\,)$          & $\{(\}$                                            & $\{(\}$ \\
16 & $F \rightarrow \text{FUNC}\,(\,E\,)$ & $\{\text{FUNC}\}$                              & $\{\text{FUNC}\}$ \\
17 & $F \rightarrow \text{NUM}$       & $\{\text{NUM}\}$                                   & $\{\text{NUM}\}$ \\
18 & $F \rightarrow \text{VAR}$       & $\{\text{VAR}\}$                                   & $\{\text{VAR}\}$ \\
\bottomrule
\end{tabular}
\caption{PRIMEROS y conjuntos de SELECCIÓN para cada una de las 18 producciones.}
\end{table}

\noindent\textbf{Verificación de la condición LL(1) (disyunción por no terminal):}
\begin{itemize}
    \item \textbf{Para $E'$:} $\text{SELECT}(2) = \{+\}$, $\text{SELECT}(3) = \{-\}$, $\text{SELECT}(4) = \{\$, )\} \implies$ intersección dos a dos $= \varnothing$.
    \item \textbf{Para $T'$:} $\text{SELECT}(6) = \{*\}$, $\text{SELECT}(7) = \{/\}$, $\text{SELECT}(8) = \{+, -, \$, )\} \implies$ intersección dos a dos $= \varnothing$.
    \item \textbf{Para $U$:} $\text{SELECT}(9) = \{+\}$, $\text{SELECT}(10) = \{-\}$, $\text{SELECT}(11) = \{(, \text{FUNC}, \text{NUM}, \text{VAR}\} \implies$ intersección dos a dos $= \varnothing$.
    \item \textbf{Para $P'$:} $\text{SELECT}(13) = \{\texttt{\^{}}\}$, $\text{SELECT}(14) = \{*, /, +, -, \$, )\} \implies$ intersección $= \varnothing$.
    \item \textbf{Para $F$:} $\text{SELECT}(15) = \{(\}$, $\text{SELECT}(16) = \{\text{FUNC}\}$, $\text{SELECT}(17) = \{\text{NUM}\}$, $\text{SELECT}(18) = \{\text{VAR}\} \implies$ intersección dos a dos $= \varnothing$.
\end{itemize}
Los no terminales $E$, $T$ y $P$ poseen una única producción. En consecuencia, la gramática cumple estrictamente la definición formal de \textbf{gramática LL(1)}.

\subsection{Tabla de Análisis Sintáctico LL(1)}
A partir de los conjuntos de selección se construye directamente la tabla LL(1) sin ningún conflicto múltiple:

\begin{table}[H]
\centering
\scriptsize
\begin{tabular}{lcccccccccc}
\toprule
\textbf{NT} & \texttt{+} & \texttt{-} & \texttt{*} & \texttt{/} & \texttt{\^{}} & \texttt{(} & \texttt{)} & \texttt{FUNC} & \texttt{NUM/VAR} & \texttt{\$} \\
\midrule
$E$  & $TE'$ & $TE'$ & & & & $TE'$ & & $TE'$ & $TE'$ & \\
$E'$ & $+TE'$ & $-TE'$ & & & & & $\varepsilon$ & & & $\varepsilon$ \\
$T$  & $UT'$ & $UT'$ & & & & $UT'$ & & $UT'$ & $UT'$ & \\
$T'$ & $\varepsilon$ & $\varepsilon$ & $*UT'$ & $/UT'$ & & & $\varepsilon$ & & & $\varepsilon$ \\
$U$  & $+U$ & $-U$ & & & & $P$ & & $P$ & $P$ & \\
$P$  & & & & & & $FP'$ & & $FP'$ & $FP'$ & \\
$P'$ & $\varepsilon$ & $\varepsilon$ & $\varepsilon$ & $\varepsilon$ & $\texttt{\^{}}U$ & & $\varepsilon$ & & & $\varepsilon$ \\
$F$  & & & & & & $(E)$ & & $\text{FUNC}(E)$ & $\text{NUM}/\text{VAR}$ & \\
\bottomrule
\end{tabular}
\caption{Tabla LL(1) sin conflictos (cada celda contiene a lo sumo una producción).}
\end{table}

\section{Algoritmo de Reconocimiento y Evaluación Semántica}
Se implementó un \textbf{parser descendente recursivo} en la clase \texttt{ParserTrigonometrico} (archivo \texttt{trig\_parser.cpp}). Cada no terminal corresponde directamente a un método (\texttt{parseE}, \texttt{parseEPrima}, \texttt{parseT}, \texttt{parseTPrima}, \texttt{parseU}, \texttt{parseP}, \texttt{parsePPrima}, \texttt{parseF}).

El flujo de ejecución en \texttt{ServicioEvaluador} consta de tres pasos ordenados:
\begin{enumerate}
    \item \textbf{Tokenización y filtrado léxico:} \texttt{EscanerAFD} convierte la cadena de entrada en tokens y omite comentarios \texttt{//}. Se verifica que no haya errores léxicos ni tokens ajenos al lenguaje.
    \item \textbf{Validación sintáctica LL(1) y detección de variables:} se ejecuta el descenso recursivo en modo validación (\texttt{validarYDetectarVariables()}). Si la expresión tiene errores de sintaxis, se reportan inmediatamente con línea y columna antes de solicitar variables al usuario. Si es válida, se registran las variables únicas en orden de aparición.
    \item \textbf{Sustitución y evaluación directa:} se solicitan al usuario los valores reales de las variables detectadas (validando formato decimal estricto) y se ejecuta \texttt{evaluar()}, calculando directamente el valor \texttt{double} durante el descenso recursivo mediante \texttt{std::sin}, \texttt{std::cos}, \texttt{std::tan} y \texttt{std::pow} de \texttt{<cmath>}.
\end{enumerate}

\subsection{Cómo se corresponden los atributos con el Parser Recursivo}
En \texttt{trig\_parser.cpp}, cada regla atribuida se ejecuta directamente sin tabla de acciones aparte:
\begin{itemize}
    \item \textbf{Atributos Sintéticos ($\uparrow$):} son los valores \texttt{double} que \textbf{retornan} los métodos \texttt{parseE()}, \texttt{parseT()}, \texttt{parseU()}, \texttt{parseP()}, \texttt{parseF()}, \texttt{parseEPrima()}, \texttt{parseTPrima()} y \texttt{parsePPrima()} (flujo de abajo hacia arriba en el árbol).
    \item \textbf{Atributos Heredados ($\downarrow$):} son los \textbf{parámetros} \texttt{acumulado} y \texttt{base} que \texttt{parseE} pasa a \texttt{parseEPrima(acum)}, \texttt{parseT} a \texttt{parseTPrima(acum)} y \texttt{parseP} a \texttt{parsePPrima(base)} (flujo de izquierda a derecha en el mismo nivel del árbol).
    \item \textbf{Unidad angular (Radianes / Grados):} la gramática LL(1) no cambia; el usuario puede seleccionar entre \textbf{Radianes} (por defecto) y \textbf{Grados} (\texttt{-g} o comando \texttt{grados} en modo interactivo). La conversión $\text{rad} = \text{grados} \cdot \pi / 180$ ocurre únicamente dentro de la evaluación semántica de $\langle F\rangle_p \rightarrow \text{FUNC}(\langle E\rangle_q)$, afectando solo al argumento de \texttt{sin/sen}, \texttt{cos} y \texttt{tan}.
    \item \textbf{Momento de cada operación:} en $\langle E'\rangle_{p\downarrow, q\uparrow} \rightarrow +\,\langle T\rangle_r\,\langle E'\rangle_{s\downarrow, t\uparrow}$ se calcula $s = p + r$ justo al obtener el atributo sintético $r$ de $\langle T\rangle_r$ y antes de pasar $s$ como atributo heredado al siguiente $\langle E'\rangle$; en $\langle P'\rangle_{p\downarrow, q\uparrow} \rightarrow \texttt{\^{}}\,\langle U\rangle_r$ se evalúa primero $r$ y luego $q = p^r$ (asociatividad derecha); y en $\langle F\rangle_p \rightarrow \text{FUNC}(\langle E\rangle_q)$ se evalúa $p = \text{FUNC}(q)$ al retornar $q$ de $\langle E\rangle_q$.
\end{itemize}

\subsection{Ejemplo de Derivación y Árbol con Atributos (Estilo Pizarra)}
Considérese la expresión del enunciado: \texttt{sen(3.14 / 2) + cos(0) * 2\^{}2} (usando aproximación visual de $\sin(1.57) \approx 0.99999968$ para claridad en los nodos).

\subsubsection*{Derivación a la izquierda (secuencia de producciones aplicadas)}
\[
\begin{aligned}
\langle E\rangle &\overset{1}{\Longrightarrow} \langle T\rangle\,\langle E'\rangle
\overset{5}{\Longrightarrow} \langle U\rangle\,\langle T'\rangle\,\langle E'\rangle
\overset{11}{\Longrightarrow} \langle P\rangle\,\langle T'\rangle\,\langle E'\rangle
\overset{12}{\Longrightarrow} \langle F\rangle\,\langle P'\rangle\,\langle T'\rangle\,\langle E'\rangle \\
&\overset{16}{\Longrightarrow} \text{sen}(\langle E\rangle)\,\langle P'\rangle\,\langle T'\rangle\,\langle E'\rangle
\overset{1,5,11,12,17,14}{\Longrightarrow} \text{sen}(3.14\,\langle T'\rangle)\,\langle P'\rangle\,\langle T'\rangle\,\langle E'\rangle \\
&\overset{7,11,12,17,14,8,4}{\Longrightarrow} \text{sen}(3.14 / 2)\,\langle P'\rangle\,\langle T'\rangle\,\langle E'\rangle
\overset{14,8}{\Longrightarrow} \text{sen}(3.14 / 2)\,\langle E'\rangle \\
&\overset{2}{\Longrightarrow} \text{sen}(3.14 / 2) + \langle T\rangle\,\langle E'\rangle
\overset{5,11,12,16,17,\dots}{\Longrightarrow} \text{sen}(3.14 / 2) + \cos(0)\,\langle P'\rangle\,\langle T'\rangle\,\langle E'\rangle \\
&\overset{14,6}{\Longrightarrow} \text{sen}(3.14 / 2) + \cos(0) * \langle U\rangle\,\langle T'\rangle\,\langle E'\rangle \\
&\overset{11,12,17}{\Longrightarrow} \text{sen}(3.14 / 2) + \cos(0) * 2\,\langle P'\rangle\,\langle T'\rangle\,\langle E'\rangle \\
&\overset{13,11,12,17,14}{\Longrightarrow} \text{sen}(3.14 / 2) + \cos(0) * 2\texttt{\^{}}2\,\langle T'\rangle\,\langle E'\rangle
\overset{8,4}{\Longrightarrow} \text{sen}(3.14 / 2) + \cos(0) * 2\texttt{\^{}}2
\end{aligned}
\]

\subsubsection*{Árbol de derivación con flujo de atributos sintéticos ($\uparrow$) y heredados ($\downarrow$)}
La Figura~1 muestra el árbol de derivación completo anotado al estilo de clase: las flechas \textbf{continuas} muestran la derivación sintáctica, las flechas \textbf{azules} indican el paso de atributos \textbf{sintéticos ($\uparrow$, de abajo hacia arriba)} y las flechas \textbf{rojas punteadas} indican el paso de atributos \textbf{heredados ($\downarrow$, de izquierda a derecha o de arriba hacia abajo)}:

\begin{figure}[H]
\centering
\begin{tikzpicture}[
  scale=0.71,
  transform shape,
  nt/.style={rectangle,rounded corners=2pt,draw=black!75,fill=blue!4,align=center,font=\scriptsize,inner sep=3pt},
  tok/.style={font=\scriptsize\ttfamily,inner sep=2pt},
  der/.style={-,draw=black!70,thick},
  syn/.style={->,>=Stealth,thick,draw=blue!75!black},
  inh/.style={->,>=Stealth,dashed,thick,draw=red!75!black}
]
% Raiz
\node[nt] (E) {$\langle E\rangle_{4.99999968}$};

% Nivel 1: T y E'
\node[nt, below left=1.1cm and 3.1cm of E] (T1) {$\langle T\rangle_{0.99999968}$};
\node[nt, below right=1.1cm and 2.8cm of E] (Ep1) {$\langle E'\rangle_{0.99999968\downarrow,\,4.99999968\uparrow}$};

\draw[der] (E) -- (T1);
\draw[der] (E) -- (Ep1);
\draw[inh] (T1) to[bend left=14] node[above,font=\tiny\color{red!75!black}]{$t=q$} (Ep1);
\draw[syn] (Ep1) to[bend right=18] node[right,font=\tiny\color{blue!75!black}]{$p=r$} (E);

% Subarbol izquierdo: sen(3.14 / 2)
\node[nt, below left=1.0cm and 0.8cm of T1] (F1) {$\langle F\rangle_{0.99999968}$};
\node[nt, below right=1.0cm and 0.4cm of T1] (Tp1) {$\langle T'\rangle_{0.99999968\downarrow,\,0.99999968\uparrow}$};
\node[tok, below=0.75cm of Tp1] (lam1) {$\lambda$};
\draw[der] (T1) -- (F1);
\draw[der] (T1) -- (Tp1);
\draw[der] (Tp1) -- (lam1);
\draw[inh] (F1) to[bend left=12] node[above,font=\tiny\color{red!75!black}]{$t=q$} (Tp1);
\draw[syn] (Tp1) to[bend right=15] node[right,font=\tiny\color{blue!75!black}]{$\uparrow$} (T1);

\node[tok, below left=0.85cm and 0.9cm of F1] (senTok) {sen(};
\node[nt, below=0.85cm of F1] (Earg) {$\langle E\rangle_{1.57}$};
\node[tok, below right=0.85cm and 0.9cm of F1] (rpar1) {)};
\draw[der] (F1) -- (senTok);
\draw[der] (F1) -- (Earg);
\draw[der] (F1) -- (rpar1);
\draw[syn] (Earg) to[bend right=26] node[right,font=\tiny\color{blue!75!black}]{$\sin(1.57)$} (F1);

\node[nt, below left=0.95cm and 0.5cm of Earg] (F314) {$\langle F\rangle_{3.14}$};
\node[nt, below right=0.95cm and 0.4cm of Earg] (TpDiv) {$\langle T'\rangle_{3.14\downarrow,\,1.57\uparrow}$};
\node[tok, below=0.7cm of F314] (n314) {3.14$_{3.14}$};
\draw[der] (Earg) -- (F314);
\draw[der] (Earg) -- (TpDiv);
\draw[der] (F314) -- (n314);
\draw[syn] (n314) to[bend left=18] (F314);
\draw[inh] (F314) to[bend left=12] (TpDiv);
\draw[syn] (TpDiv) to[bend right=18] (Earg);

\node[tok, below left=0.85cm and 0.7cm of TpDiv] (divTok) {/};
\node[nt, below=0.85cm of TpDiv] (F2a) {$\langle U\rangle_{2}$};
\node[nt, below right=0.85cm and 0.5cm of TpDiv] (TpDivEnd) {$\langle T'\rangle_{1.57\downarrow,\,1.57\uparrow}$};
\node[tok, below=0.65cm of F2a] (n2a) {2$_{2}$};
\node[tok, below=0.65cm of TpDivEnd] (lamDiv) {$\lambda$};
\draw[der] (TpDiv) -- (divTok);
\draw[der] (TpDiv) -- (F2a);
\draw[der] (TpDiv) -- (TpDivEnd);
\draw[der] (F2a) -- (n2a);
\draw[der] (TpDivEnd) -- (lamDiv);
\draw[syn] (n2a) to[bend left=18] (F2a);
\draw[inh] (F2a) to[bend left=10] node[above,font=\tiny\color{red!75!black}]{$s=3.14/2$} (TpDivEnd);
\draw[syn] (TpDivEnd) to[bend right=18] (TpDiv);

% Subarbol derecho: + cos(0) * 2^2
\node[tok, below left=1.0cm and 1.8cm of Ep1] (plusTok) {+};
\node[nt, below=1.0cm of Ep1] (T2) {$\langle T\rangle_{4.0}$};
\node[nt, below right=1.0cm and 1.9cm of Ep1] (Ep2) {$\langle E'\rangle_{4.99999968\downarrow,\,4.99999968\uparrow}$};
\node[tok, below=0.75cm of Ep2] (lamEnd) {$\lambda$};
\draw[der] (Ep1) -- (plusTok);
\draw[der] (Ep1) -- (T2);
\draw[der] (Ep1) -- (Ep2);
\draw[der] (Ep2) -- (lamEnd);
\draw[inh] (T2) to[bend left=12] node[above,font=\tiny\color{red!75!black}]{$s=p+r$} (Ep2);
\draw[syn] (Ep2) to[bend right=18] node[right,font=\tiny\color{blue!75!black}]{$q=t$} (Ep1);

\node[nt, below left=1.0cm and 0.8cm of T2] (Fcos) {$\langle F\rangle_{1.0}$};
\node[nt, below right=1.0cm and 0.6cm of T2] (Tp2) {$\langle T'\rangle_{1.0\downarrow,\,4.0\uparrow}$};
\draw[der] (T2) -- (Fcos);
\draw[der] (T2) -- (Tp2);
\draw[inh] (Fcos) to[bend left=12] node[above,font=\tiny\color{red!75!black}]{$t=q$} (Tp2);
\draw[syn] (Tp2) to[bend right=18] (T2);

\node[tok, below=0.75cm of Fcos] (cos0) {cos(0$_{0}$)};
\draw[der] (Fcos) -- (cos0);
\draw[syn] (cos0) to[bend left=18] node[left,font=\tiny\color{blue!75!black}]{$\cos(0)$} (Fcos);

\node[tok, below left=0.95cm and 0.9cm of Tp2] (mulTok) {*};
\node[nt, below=0.95cm of Tp2] (Ppow) {$\langle P\rangle_{4.0}$};
\node[nt, below right=0.95cm and 0.9cm of Tp2] (Tp2End) {$\langle T'\rangle_{4.0\downarrow,\,4.0\uparrow}$};
\node[tok, below=0.65cm of Tp2End] (lamMul) {$\lambda$};
\draw[der] (Tp2) -- (mulTok);
\draw[der] (Tp2) -- (Ppow);
\draw[der] (Tp2) -- (Tp2End);
\draw[der] (Tp2End) -- (lamMul);
\draw[inh] (Ppow) to[bend left=10] node[above,font=\tiny\color{red!75!black}]{$s=1*4$} (Tp2End);
\draw[syn] (Tp2End) to[bend right=18] (Tp2);

\node[nt, below left=0.85cm and 0.3cm of Ppow] (F2b) {$\langle F\rangle_{2}$};
\node[nt, below right=0.85cm and 0.3cm of Ppow] (Pp) {$\langle P'\rangle_{2\downarrow,\,4\uparrow}$};
\node[tok, below=0.65cm of F2b] (n2b) {2$_{2}$};
\node[tok, below=0.65cm of Pp] (pow2) {\texttt{\^{}} 2$_{2}$};
\draw[der] (Ppow) -- (F2b);
\draw[der] (Ppow) -- (Pp);
\draw[der] (F2b) -- (n2b);
\draw[der] (Pp) -- (pow2);
\draw[syn] (n2b) to[bend left=18] (F2b);
\draw[inh] (F2b) to[bend left=10] (Pp);
\draw[syn] (pow2) to[bend right=18] node[right,font=\tiny\color{blue!75!black}]{$2^2$} (Pp);
\draw[syn] (Pp) to[bend right=16] (Ppow);
\end{tikzpicture}
\caption{Árbol de derivación con atributos para \texttt{sen(3.14 / 2) + cos(0) * 2\^{}2}. Las flechas azules muestran el flujo sintético ($\uparrow$, abajo hacia arriba) y las rojas punteadas el flujo heredado ($\downarrow$, izquierda a derecha).}
\end{figure}

\subsection*{Validaciones Implementadas}
\begin{itemize}
    \item \textbf{Léxicas:} caracteres no reconocidos (\texttt{@}), tokens de C no permitidos (\texttt{;}, cadenas), literales hexadecimales (\texttt{0x10}) y puntos aislados (\texttt{.}, \texttt{3..1}). Se admiten enteros, decimales (\texttt{3.14}, \texttt{.5}, \texttt{3.}) y notación científica (\texttt{1.5e-3}, \texttt{2E3}).
    \item \textbf{Sintácticas:} entrada vacía, paréntesis vacíos o desbalanceados, funciones sin paréntesis (\texttt{sin x}), funciones no soportadas (\texttt{log(x)}), exponentes incompletos o letras pegadas (\texttt{3445ea}, \texttt{1.5e+}, \texttt{1.2.3}), operadores consecutivos o faltantes (\texttt{2 * / 3}, \texttt{2 x}) y límite de recursión (512 niveles).
    \item \textbf{Matemáticas:} división por cero (\texttt{divisor == 0}), \texttt{0} elevado a exponente negativo, base negativa con exponente fraccionario (\texttt{(-4)\^0.5}), constantes fuera de rango (\texttt{1e400}) y desbordamientos (\texttt{!std::isfinite}).
    \item \textbf{Entrada de variables:} rechazo de cadenas vacías, texto no numérico, \texttt{nan}, \texttt{inf}, \texttt{3445ea} o \texttt{1.5e+} (acepta enteros, decimales y notación científica válida).
\end{itemize}

\section{Casos de Prueba y Validación Matemática}
Se ejecutaron 53 verificaciones unitarias automatizadas (\texttt{make test}) y se contrastaron todas las expresiones del enunciado, distintos formatos de constantes reales y funciones trigonométricas anidadas contra \texttt{math} de Python 3 (con $x=1.0$, $y=0.5$, $z=2.0$ en radianes):

\begin{table}[H]
\centering
\scriptsize
\begin{tabular}{p{5.8cm}rrr}
\toprule
\textbf{Expresión de entrada} & \textbf{Evaluador C++} & \textbf{Python \texttt{math}} & \textbf{$|\Delta|$} \\
\midrule
\texttt{sen(3.14 / 2) + cos(0) * 2\^{}2} & \texttt{4.99999968293183} & \texttt{4.99999968293183} & $4.44\times 10^{-15}$ \\
\texttt{sen(x) + cos(y)} & \texttt{1.71905354669827} & \texttt{1.71905354669827} & $6.66\times 10^{-16}$ \\
\texttt{sen(cos(x))+ tan(y)/cos(x)} & \texttt{1.52550049333466} & \texttt{1.52550049333466} & $6.66\times 10^{-16}$ \\
\texttt{cos(x)\^{}2 + sen(x)\^{}2} & \texttt{1} & \texttt{1} & $0.00$ \\
\texttt{sin(x)} & \texttt{0.841470984807897} & \texttt{0.841470984807897} & $4.44\times 10^{-16}$ \\
\texttt{cos(x + y)} & \texttt{0.0707372016677029} & \texttt{0.0707372016677029} & $0.00$ \\
\texttt{tan(x\^{}2 + sin(x))} & \texttt{-3.60380353200881} & \texttt{-3.60380353200881} & $4.44\times 10^{-15}$ \\
\texttt{.5 + 3.} & \texttt{3.5} & \texttt{3.5} & $0.00$ \\
\texttt{-.25 + +.75} & \texttt{0.5} & \texttt{0.5} & $0.00$ \\
\texttt{1.5e-3 * 2E3} & \texttt{3} & \texttt{3} & $0.00$ \\
\texttt{.5e2 + 3.e-1} & \texttt{50.3} & \texttt{50.3} & $0.00$ \\
\texttt{sen(cos(tan(0.5)))} & \texttt{0.754210756032613} & \texttt{0.754210756032613} & $4.44\times 10^{-16}$ \\
\texttt{tan(sin(cos(.5)) + 1.2e-2)} & \texttt{0.991631496487538} & \texttt{0.991631496487538} & $1.11\times 10^{-16}$ \\
\texttt{cos(sen(x) + tan(y * 1e-1))} & \texttt{0.628235835763506} & \texttt{0.628235835763506} & $1.11\times 10^{-16}$ \\
\texttt{-2\^{}2 + 10 / 2 * 3} & \texttt{11} & \texttt{11} & $0.00$ \\
\texttt{2\^{}3\^{}2} & \texttt{512} & \texttt{512} & $0.00$ \\
\texttt{sin(-x + 0.5 * z) / (1 + cos(y)\^{}2)} & \texttt{0} & \texttt{0} & $0.00$ \\
\bottomrule
\end{tabular}
\caption{Comparación de resultados obtenidos vs. referencia matemática (\texttt{x=1.0, y=0.5, z=2.0}).}
\end{table}

Asimismo, los 19 casos inválidos contenidos en \texttt{tests/expresiones\_invalidas.txt} fueron rechazados en su fase correspondiente (léxica, sintáctica o matemática) con indicación precisa de línea, columna y lexema problemático (ver \texttt{output/resultados\_pruebas.txt}).

\subsection*{Ejecución Interactiva e Ingreso de Valores de Variables}
Cuando la expresión contiene variables simbólicas (ej. \texttt{x}, \texttt{y}), al ejecutar \texttt{./evaluador} (o \texttt{./evaluador -e "<expr>"} sin \texttt{-v}), el módulo valida primero la sintaxis LL(1), detecta las variables únicas en orden de aparición y las solicita una a una al usuario:

\begin{lstlisting}
$ ./evaluador
Seleccione unidad angular [1 = Radianes (por defecto), 2 = Grados]: 1
[RAD] > sen(cos(x))+ tan(y)/cos(x)
  Ingrese valor para 'x' (radianes si es angulo): abc
  Valor invalido. Ingrese un numero real (ej: 0, 3.14, -1.5, .5, 1.5e-3).
  Ingrese valor para 'x' (radianes si es angulo): 1.0
  Ingrese valor para 'y' (radianes si es angulo): 0.5
Resultado: 1.52550049333466

[RAD] > grados
Modo cambiado a GRADOS.
[DEG] > sen(a + c)
  Ingrese valor para 'a' (grados si es angulo): 45
  Ingrese valor para 'c' (grados si es angulo): 45
Resultado: 1
\end{lstlisting}

También es posible pasar valores y unidad directamente desde la línea de comandos mediante \texttt{-v x=1.0 -v y=0.5} y \texttt{-g} (grados).

\section{Conclusiones}
La gramática LL(1) diseñada resuelve de forma natural la precedencia y asociatividad de operadores aritméticos, unarios, potencias y funciones trigonométricas sin ambigüedad. El uso de la arquitectura modular de \texttt{scanner\_c} permitió construir el evaluador de forma limpia y mantenible, garantizando una detección temprana y precisa de errores en todas las fases del análisis.

\end{document}
